% Generated draft: requires mathematical and source review.
\documentclass[11pt]{article}
\usepackage[margin=25mm]{geometry}
\usepackage{fontspec}
\setmainfont{Latin Modern Roman}
\usepackage{amsmath,amssymb,mathtools,unicode-math}
\setmathfont{Latin Modern Math}
\usepackage{enumitem,needspace,xurl,hyperref}
\setlength{\parindent}{0pt}
\setlength{\parskip}{5pt}
\setlength{\emergencystretch}{3em}
\newcommand{\N}{\mathbb{N}}
\newcommand{\Z}{\mathbb{Z}}
\newcommand{\Q}{\mathbb{Q}}
\newcommand{\R}{\mathbb{R}}
\newcommand{\C}{\mathbb{C}}
\newcommand{\F}{\mathbb{F}}
\newcommand{\A}{\mathbb{A}}
\newcommand{\PP}{\mathbb{P}}
\newcommand{\E}{\mathbb{E}}
\newcommand{\eps}{\varepsilon}
\newcommand{\dd}{\,\mathrm{d}}
\begin{document}
\section*{UnsolvedMath ID: 30003338}
\textbf{Source problem code:} catalogue entry
\subsection{Definitions and mathematical statement}
Let \(G=(V,E)\) be a finite bipartite graph with specified parts \(A\) and \(B\), so that \(V=A\sqcup B\) and every edge has one endpoint in each part. Let \(q\in\N\) satisfy \(q\geq3\), and use the color set \(C_q=\{0,1,\ldots,q-1\}\). Define the finite set of proper colorings by
\[
\Omega_q(G)=\{f:V\to C_q: f(u)\neq f(v)\text{ for every }uv\in E\}.
\]
This set is nonempty because \(G\) is bipartite. Choose \(f\) uniformly from \(\Omega_q(G)\), and define
\[
S=f^{-1}(\{0\})\cap A,\qquad \eta=(\mathbf{1}_{\{a\in S\}})_{a\in A}\in\{0,1\}^A.
\]
Equip \(\{0,1\}^A\) with the coordinatewise partial order. A function \(X:\{0,1\}^A\to\R\) is increasing if \(X(x)\leq X(y)\) whenever \(x_a\leq y_a\) for every \(a\in A\). The random set \(S\) has positive association if, for every pair of increasing functions \(X,Y:\{0,1\}^A\to\R\),
\[
\E[X(\eta)Y(\eta)]\geq\E[X(\eta)]\,\E[Y(\eta)],
\]
where expectations are taken over the uniform proper coloring.

The question is whether this inequality holds for every finite bipartite graph with specified parts \(A,B\), every integer \(q\geq3\), and every such pair \(X,Y\). The supplied record labels this question open. Its background also mentions an unresolved four-vertex test, but does not specify that test; no additional formulation of it is asserted here.
\subsection{Short English statement}
For any finite bipartite graph and any integer \(q\geq3\), choose a proper \(q\)-coloring uniformly. Must any two increasing real-valued functions of the set of vertices in \(A\) colored \(0\) have nonnegative covariance? The supplied open-status label and verification claims require independent human review.
\subsection{Sources}
\begin{itemize}
\item Local UnsolvedMath registry, permanent ID \(30003338\), title ``Positive Association in Random Proper Colorings,'' source problem code \texttt{catalogue entry}; collection: Oberwolfach Reports Open Problems. The record gives proposed year \(2017\) and names no proposer.
\item Supplied source citation: \emph{Combinatorics (2018)}, Oberwolfach Reports, DOI \texttt{10.4171/owr/2017/1}, pp.~5--81, \url{https://doi.org/10.4171/owr/2017/1}. The supplied literature-sources entry labels this same link ``original OWR report.'' Oberwolfach Reports is described in the record as documenting mathematical workshops and their open-problem discussions.
\item Supplied literature assessment, checked \texttt{2026-08-21}: the question is labeled open; the assessment says that the source report records the positive-association inequality as unknown, even for a four-vertex test, and that no later primary resolution of this exact question was located. The listed source checked is the original OWR report at \url{https://doi.org/10.4171/owr/2017/1}. This is supplied dated curation, not an independently verified literature assessment.
\item The record carries statement status \texttt{corrected\_verified} and reports that its clean statement was checked against the official Oberwolfach report and preserves the report's hypotheses, notation, and target. These are supplied verification claims; the source has not been independently consulted for this draft. The supplied citation's title year \(2018\), DOI year component \(2017\), and proposed year \(2017\) are retained without resolving their relationship.
\end{itemize}
\end{document}
